前三章给出的是定义与定理。定义本身不会告诉你什么时候该把它拿出来用，本章处理的共同困难就是这一步：把一个初看不像随机过程的问题，改写成一个随机过程，然后让已有的定理替你回答。

这一步之所以值得单独讲，是因为改写的方式决定了你能问什么。同一个赌徒破产问题，写成 Markov 链就该用首步分析逐层递推，写成 Martingale 就该找守恒量一步跨过去；两条路给出同一个 \(i/N\)，但前者顺带算出了期望时间 \(i(N-i)\)，后者只用到``公平''这一个事实，因而在完全不知道链的结构时仍然可用。选哪一条不是品味问题，是你手上有什么、想要什么。

四篇各演示一个建模动作，彼此不重叠：把独立小扰动累加起来看整体尺度；把``没有可预测漂移''单独提出来，再追问自选停手时刻能不能占到便宜；把``重要性''定义成一条人造随机游走的平稳分布；在状态不可见时，用观测和动力学一起把状态反推回来。

\begin{itemize}
\tightlist
\item \hyperref[sec:06-Random-Processes/04-Applications/01]{``随机游走：从局部步进到整体尺度''}——步数与位移之间隔着一个平方根，这条标度决定了要多少数据才看得出趋势；
\item \hyperref[sec:06-Random-Processes/04-Applications/02]{``Martingale：公平游戏与停时''}——公平游戏无法被停止策略撬开，条件一旦失效则相反，加倍下注与``看到显著就停''是同一个漏洞；
\item \hyperref[sec:06-Random-Processes/04-Applications/03]{``PageRank：把网页排序变成平稳概率''}——排序问题变成求平稳分布，阻尼因子买来的是不可约与非周期；
\item \hyperref[sec:06-Random-Processes/04-Applications/04]{``Hidden Markov Model：从观测反推状态''}——条件独立结构让指数条路径的求和塌缩成逐列递推。
\end{itemize}

按依赖关系，前两篇是主线：随机游走的 \(\sqrt n\) 标度在后面的扩散与随机微分方程里还要用两次，Martingale 的停时条件则直接迁移到实验分析和序贯决策。后两篇更接近独立的案例，可以按兴趣挑读；若你的工作与序列建模有关，第四篇是通向 Kalman 滤波和状态空间模型的入口。

读的时候有四个问题可以一路带着：状态里装的信息够不够，局部转移怎样累积成长期行为，新到的信息改变了哪个条件期望或后验，以及算法输出的含义究竟由哪些建模选择定义。最后一个尤其容易被跳过——PageRank 那一篇会说明，一个排名分数的全部意义都写在你选的那个随机浏览模型里。
